% Appendix: side channels in detail. Moved from Section~\ref{sec:side}. Sources:
% benchmarks/paper/PREREG_X3_X6.md on branch paper/x3-x6 (X6, with amendments 1 to 3, all made before
% any full result); results from evidence/x3-x6-run1.md (GitHub Actions run 40940385704, commit
% a191626). Every number is as x6_stats.py printed it; the two places where the script departs from
% the written rules are stated in the text. The follow-up analyses (amendment 4, exploratory,
% specified in writing after the main results and before running; run 40982986417, commit 9e22b89)
% are from evidence/x3-x6-followup.md, as x3x6_followup.py printed them; they change no verdict.
% "Reveal" is used throughout; "leak" is reserved for returned content.

\section{Side channels in detail (X6)}\label{app:side}

This appendix gives the setup, the hypotheses and the full results of the side-channel experiment, including the
forensics of the build offsets, which Section~\ref{sec:side} summarizes. Beside the game's differencing attacker
(Section~\ref{sec:side-game}), a forced-choice attacker whose advantage is $2\,\mathrm{AUC} - 1$ and for whom $|\mathrm{AUC} - 0.5|$ is half the
size of that advantage, we report a \emph{one-shot attacker}, who sees one observation and must threshold it; its
success is the AUC over all $W_1$ against all $W_0$ observations. The pre-registration also asks for the one-shot attacker's largest true-positive rate minus false-positive rate over thresholds; it is in the released data, as \texttt{ks} in \texttt{x6\_cells.csv}. For post-filter latency at $m = 100$ it was 0.64 to 0.66 at $p = 10\%$ and 0.88 to 0.93 at $p = 50\%$.

\subsection{Setup}\label{sec:side-setup}

\paragraph{Data and probes.}
The background is all 57,638 FiQA passages~\cite{maia2018www18}, with its 6,648 queries as probe
candidates; far chunks come from 12,000 SciDocs abstracts, both from BEIR~\cite{thakur2021beir}, and
all text is embedded with all-MiniLM-L6-v2 (384 dimensions)~\cite{reimers2019sentencebert,wang2020minilm}.
Each of five seeds places eight groups of nine probes, one with $m = 1{,}000$ secrets, two with 100,
three with 10 and three with 1. A probe's secrets are its $m$ nearest passages, excluding hubs, the 1\% of
passages found among the 2,000 nearest neighbors of the largest number of queries; within a group, no probe's secrets fall in the
zone a group mate can explore, its nearest passages up to 1,000 beyond its own secrets, which is the
post-filter's cap. One index $W_1$ holds every group's secrets; for each group $g$, an index $W_0^g$
replaces only group $g$'s 1,233 secrets by far chunks. A probe of group $g$ thus compares two worlds
that differ in its own secrets and in those of its eight group mates, which the zone rule places
beyond the depth the layer explores for it; collection statistics, however, count them all.
Background passages fall into 100 permission buckets, and three
attackers read 1\%, 10\% and 50\% of them; secrets carry an ACL none of them holds. Every seed
builds every index afresh, $W_1$ first, so every $W_0^g$ is built later in its seed than $W_1$ and build age is aligned with world in every latency cell.

\paragraph{Configurations.}
The \emph{pre-filter} runs each store's approximate filtered search at the settings the pre-registration fixed, which it calls documented defaults:
pgvector (HNSW $m = 16$, \texttt{ef\_construction} 64, \texttt{ef\_search} 40, with iterative scans in relaxed order, which pgvector's own default leaves off), Qdrant (default HNSW) and Elasticsearch (\texttt{num\_candidates} 100, against its default of 15 at $k = 10$), with FAISS over a flat index and \sys{}'s in-memory backend as
exact controls. The \emph{post-filter} is \sys{}'s widening loop at its defaults: fetch $4k$, widen
fourfold until $k$ hits are readable, the store runs out, or 1,000 candidates have been examined.
Four mitigations are measured against them. \emph{Constant-work widening} (M1) fetches the cap of
1,000 in one round, so every request examines the same number of candidates. \emph{Exact
pre-filtered search} (M2) ranks every readable chunk: pgvector with index scans disabled, Qdrant's
exact flag, Elasticsearch's \texttt{script\_score}. \emph{Separation} (M3) gives the attacker an index
that holds only its own tenant's chunks, as a tenant-per-index layout would; its advantage is zero by
construction, and measuring it twice provides an A/A null for the harness. \emph{Latency padding}
(M4) pads each answer to the next multiple of a budget fixed on the first seed and is evaluated on
the other four. No padder ran: the padded latency is computed afterwards from the recorded latencies, so it models an ideal padder and ignores scheduler jitter. For collection statistics,
keyword search runs under the pre-filter on Elasticsearch (BM25), on Qdrant (sparse vectors with an
IDF modifier) and on pgvector (\texttt{ts\_rank\_cd}, which uses no collection statistics: a
negative control), and Elasticsearch adds reciprocal-rank fusion of vector and keyword search.

\paragraph{Measurement and statistics.}
For every seed, store, probe, attacker and configuration, one warm-up is followed by five timed
repetitions, each of which runs $W_0^g$, $W_1$ and the two separated indexes in a fresh random order,
so drift affects every world alike; an observation is the median over repetitions. The in-memory
control, which rescores the whole corpus on every search, runs a reduced design fixed in advance:
three probes per level and seed, $p = 10\%$ only, three repetitions. Every retrieval is
re-checked, and a hit outside the attacker's readable set would stop the run. Intervals are 95\%
cluster-bootstrap intervals over (seed, probe) pairs with 10,000 resamples; they do not resample index builds, the demonstrated source of latency noise, so latency intervals are too narrow. A channel is
\emph{present} when the paired AUC reaches the hypothesis's threshold and its interval's lower end
exceeds 0.55, and \emph{absent} when the interval lies within $[0.40, 0.60]$. Two guards precede any
verdict: the A/A comparison of the two separated indexes must show no latency channel, and exact
configurations must return identical hit lists in $W_0^g$ and $W_1$ for at least 99\% of probes. The
experiment runs on GitHub-hosted runners with four virtual CPUs and 16\,GB of memory.

\begin{table}[tbp]
\centering
\small
\begin{tabular}{@{}p{0.20\linewidth}p{0.42\linewidth}p{0.32\linewidth}@{}}
\toprule
Hypothesis & Prediction & Verdict \\
\midrule
H1 widening & Post-filter: count AUC $\ge 0.95$ at $m = 1{,}000$; latency AUC $\ge 0.75$ at $m = 100$
  ($p = 10\%, 50\%$); no channel at $m = 1$ & Falsified: count 0.613 to 0.688 at $p = 50\%$;
  FAISS latency 0.667 at $m = 1$ \\
H2 constant work (M1) & No latency channel; the count channel remains at $m = 1{,}000$ & Falsified:
  count as in H1; a reverse latency channel on pgvector \\
H3 approximate pre-filter & pgvector and Elasticsearch: latency AUC $\ge 0.75$ at $m = 1{,}000$;
  Qdrant: no channel at $p = 1\%, 10\%$ & pgvector confirmed; Elasticsearch inconclusive; Qdrant
  confirmed \\
H4 exact pre-filter (M2) & No channel in any observable & Inconclusive: identical answers; through latency, build offsets on pgvector and FAISS, wide intervals at 40 clusters on Qdrant and Elasticsearch; the in-memory control failed G1 \\
H5 collection statistics & BM25 and IDF-weighted sparse: score AUC $\ge 0.90$ at every $m$;
  \texttt{ts\_rank\_cd}: none; fusion changes the hit list for at least 25\% of probes at
  $m = 1{,}000$ & BM25, sparse: inconclusive (0.794 to 0.886 at $m \le 10$);
  \texttt{ts\_rank\_cd}, fusion: confirmed \\
H6 padding (M4) & No latency channel wherever the unpadded AUC was $\ge 0.75$ (forward channels
  only) & Inconclusive as printed (script defect); confirmed under the written rule, by construction \\
\bottomrule
\end{tabular}
\caption{Pre-registered hypotheses of the side-channel experiment, abbreviated, with the verdicts of
the full run (Section~\ref{sec:side-results}); the full rules, including what falsifies each, are in
the released pre-registration.}
\label{tab:side-hyp}
\end{table}

\subsection{Full results}\label{app:side-results}

\begin{table}[tbp]
\centering
\small
\begin{tabular}{@{}llrrrr@{}}
\toprule
Configuration & Observable & $m = 1$ & 10 & 100 & 1,000 \\
\midrule
Post-filter, FAISS & latency & 0.67 & 0.67 & 0.80 & 0.76 \\
Post-filter, pgvector & latency & 0.37 & 0.38 & 0.57 & 0.68 \\
Post-filter, Qdrant & latency & 0.47 & 0.46 & 0.75 & 0.78 \\
Post-filter, Elasticsearch & latency & 0.48 & 0.58 & 0.76 & 0.86 \\
Post-filter, all four & count & 0.50--0.51 & 0.50--0.51 & 0.58--0.59 & 0.84--0.87 \\
\midrule
Constant work, FAISS & latency & 0.63 & 0.57 & 0.46 & 0.13 \\
Constant work, pgvector & latency & 0.29 & 0.22 & 0.07 & 0.02 \\
Constant work, Qdrant & latency & 0.47 & 0.47 & 0.45 & 0.35 \\
Constant work, Elasticsearch & latency & 0.57 & 0.57 & 0.55 & 0.62 \\
\midrule
Pre-filter, pgvector & latency & 0.33 & 0.44 & 0.79 & 0.88 \\
Pre-filter, pgvector & recall & 0.49 & 0.47 & 0.31 & 0.22 \\
Pre-filter, Qdrant & latency & 0.47 & 0.43 & 0.47 & 0.57 \\
Pre-filter, Elasticsearch & latency & 0.49 & 0.50 & 0.48 & 0.60 \\
\midrule
Exact, FAISS (pre-filter) & latency & 0.25 & 0.30 & 0.25 & 0.35 \\
Exact, pgvector & latency & 0.17 & 0.14 & 0.13 & 0.12 \\
Exact, Qdrant & latency & 0.45 & 0.46 & 0.50 & 0.42 \\
Exact, Elasticsearch & latency & 0.50 & 0.48 & 0.47 & 0.47 \\
Exact, all four & count, recall, score & 0.50 & 0.50 & 0.50 & 0.50 \\
\midrule
BM25, Elasticsearch & score & 0.79 & 0.88 & 0.98 & 0.99 \\
Sparse IDF, Qdrant & score & 0.81 & 0.89 & 0.98 & 1.00 \\
\texttt{ts\_rank\_cd}, pgvector & score & 0.50 & 0.50 & 0.50 & 0.50 \\
\texttt{ts\_rank\_cd}, pgvector & latency & 0.45 & 0.47 & 0.50 & 0.70 \\
\bottomrule
\end{tabular}
\caption{Paired AUC of the differencing attacker, $W_1$ against $W_0^g$, pooled over the three
attackers: 120, 120, 80 and 40 probes for $m = 1$, 10, 100 and 1,000 over five seeds. Above 0.5 the
observable moves in the direction predicted for $W_1$ (latency up; count, recall and score down);
below 0.5 it moves the other way, which serves an attacker who knows the sign just as well, so the
size of a channel is $|\mathrm{AUC} - 0.5|$. Latency values must be read against the differences between
builds of identical content, which reached, per $m$, $|\mathrm{AUC} - 0.5|$ of 0.16 for FAISS's post-filter, 0.38
for its pre-filter and 0.33 for pgvector's exact search, and 0.08 for Qdrant's exact search
(Section~\ref{sec:side-results}); the other configurations were not measured that way. No correction
for multiple comparisons is applied: the verdicts rest on the cells the hypotheses name, and the rest
of the table is descriptive. Intervals
for the cells the hypotheses name are in the text; all are in the released report.}
\label{tab:side-auc}
\end{table}

\paragraph{The run.}
Every pre-registered cell ran: five seeds for each store, all 72 probes placed in every seed, and
520,320 retrievals without an error. None returned a chunk outside the attacker's readable set. In
three seeds the SciDocs passages far enough from every probe (2,331 to 4,864) were fewer than the
secrets (8,417 to 8,985), and the harness reused far chunks rather than drawing them without
replacement as the pre-registration states; every far chunk still lay outside every probe's zone.
Table~\ref{tab:side-auc} summarizes the channels by $m$. Verdicts are as the analysis script printed
them; where the script departs from the written rules, we say so.

\paragraph{Guards.}
Exact search returned identical hit lists in $W_1$ and $W_0^g$ for every probe and attacker (G2:
change rate 0.000 for pgvector, Qdrant and Elasticsearch with exact search and for FAISS). The A/A
guard G1 held for all 20 configurations of the four vector stores, with paired AUCs from 0.463 (95\%
interval 0.436 to 0.492) to 0.536 (0.507 to 0.563) over 360 clusters. It failed for the in-memory
control (0.533, interval 0.400 to 0.650; and 0.500, 0.367 to 0.633): with 60 clusters and three
repetitions, an interval inside $[0.40, 0.60]$ was out of reach, so its latency parts are
inconclusive, as they would have been in any case (pooled latency 0.383, 0.267 to 0.500, and 15
clusters per level). G1 compares one index
with itself; it cannot see a difference between two builds, and such differences appear below.

\paragraph{Widening (H1, falsified).}
Section~\ref{sec:side-results} gives the paired AUCs at $m = 100$. There FAISS took 14.8 against 7.2\,ms on average at $p = 50\%$, and the
one-shot attacker did almost as well (pooled AUC 0.962 for FAISS at $p = 50\%$). pgvector's 0.662
at $p = 10\%$ was inconclusive, although $W_1$ took 18.7 against 10.7\,ms on average:
on probes that needed no extra round, its post-filter was slightly faster on $W_1$, as it was at
$m = 1$ (0.369) and $m = 10$ (0.383). At
$m = 1{,}000$ an attacker reading 1\% of the background received a short answer in $W_1$ for every
probe on every store (count AUC 1.000).

Two parts failed. At $m = 1{,}000$ and $p = 50\%$ the count AUC was 0.613 to 0.688, below the
falsification bound of 0.80 on every store, and at $p = 10\%$ it was 0.912 to 0.925. The prediction
assumed that the 1,000 nearest candidates are all secrets. Hubs may not be secrets, so they remain in
the background, interleaved with the secrets: a probe's 1,000th secret ranked as deep as 1,040 to
1,271 (the deepest per seed), and readable hubs among the post-filter's 1,000 candidates filled the
answer. On FAISS, $W_1$ returned 0.375, 3.35 and 7.60 hits on average at $p = 1\%$, 10\% and 50\%, and
its answer was complete for 62.5\% of probes at $p = 50\%$, which accounts for the AUC of 0.688, since
every $W_0^g$ answer was complete. A follow-up check, exploratory, specified in writing after the
main results and before running, confirmed the mechanism probe by probe: the 1,000 nearest passages of each
$m = 1{,}000$ probe held between 1 and 231 hubs (45.7 on average) and no passage that was neither a
hub nor a secret, and FAISS returned exactly the smaller of 10 and the number of readable hubs among
them, for all 40 probes and every attacker; the same count over the background predicted every
$W_0^g$ answer. The count channel is only as strong as readable content among the secrets is scarce.

The second failure is a latency difference on FAISS at $m = 1$: paired AUC 0.667 (0.617 to 0.717).
Its answers were identical in both worlds for all 360 pairs, the mean difference was 0.05 to 0.14\,ms
on 7 to 46\,ms, and the AUC was the same at $m = 10$ (0.672), with ten times as many secrets among the
candidates. Configurations whose answers are identical in both worlds show such offsets, of either
sign, at every $m$: FAISS's exact pre-filter was faster on $W_1$ (paired AUC 0.25 to 0.35), and
pgvector's exact search was 1.5 to 3.7 times slower on $W_0^g$ (0.12 to 0.17).

Two follow-up analyses, likewise exploratory and specified in writing before they ran, examined this
offset. In the
main run's records the whole difference lay in FAISS's own search call: 0.115\,ms (95\% interval
0.076 to 0.154; paired AUC 0.681) against 0.106\,ms of latency, with the same number of widening
rounds in all 360 pairs and no difference in $W_1$'s direction in the authorizer ($-0.001$\,ms) or
the rest of the layer ($-0.010$\,ms). A second run then built, for each group, $W_1$ again and
$W_0^g$ twice, just before the group, and measured them with the original $W_1$ in random order. It
reproduced the direction of the offset, more weakly ($W_1$ against the first $W_0^g$: 0.597, 0.539
to 0.656, at $m = 1$), and two builds of identical content differed by as much: $W_1$ against its
own rebuild 0.600 (0.547 to 0.653), the second build of $W_0^g$ against the first 0.361 (0.308 to
0.414), at every $m$ alike. With builds matched in time, the difference in content gave 0.447 (0.392
to 0.506) at $m = 1$, but 0.742 (0.700 to 0.783) at $m = 1{,}000$, where the widening channel is real.
The reading we fixed in advance called the offset a build difference only if that $m = 1$ interval
lay within $[0.40, 0.60]$; it reaches 0.392, so the formal reading is inconclusive. What the numbers
show is that two FAISS indexes of the same content differ in search time by as much as $W_1$ and
$W_0^g$ did, while the secret adds no detectable offset of its own at $m = 1$.

These same-content comparisons set the noise level against which every latency channel of
Table~\ref{tab:side-auc} must be read. Two builds of identical content gave paired latency AUCs of
0.342 to 0.608 for FAISS's post-filter, depending on $m$ and on which builds were compared
($|\mathrm{AUC} - 0.5|$ up to 0.16). Pooled over $m$ they were 0.257 and 0.426 for its
pre-filter (per $m$, up to 0.38, at $m = 1{,}000$), 0.186 and 0.789 for pgvector's exact search (per $m$, up to 0.33; from the tables'
vacuum state, below), and 0.494 and 0.505 for Qdrant's exact search (per $m$, up to 0.08, at $m = 1{,}000$, where its cell reads 0.42); we have no such
measurement for the other configurations, and Elasticsearch's was lost. A latency AUC between about
0.6 and 0.8 is therefore not by itself evidence about the secrets. To separate a channel from an
offset we use a criterion adopted after seeing these results: a channel grows with $m$, as in the widening loop and pgvector's iterative scan, while
an offset stays the same at every $m$, as in the configurations whose answers are identical.

\paragraph{Constant work (H2, falsified).}
Fetching the cap in one round removed the widening channel from latency pooled over $m$ and $p$ on
FAISS (0.517, 0.484 to 0.549) and Qdrant (0.456, 0.422 to 0.491), and left the count channel, which
falsifies H2 at $p = 50\%$ as it does H1. Timing did not become uninformative. Where the candidates
are secrets, $W_1$ answered \emph{faster}: on pgvector the paired AUC fell from 0.29 at $m = 1$ to 0.017
(0.000 to 0.050) at $m = 1{,}000$, where $W_1$ took 23.5 against 26.4\,ms; FAISS fell to 0.133 (0.075
to 0.200) and Qdrant to 0.350 (0.225 to 0.483), while Elasticsearch drifted the other way (0.573
pooled, 0.625 at $m = 1{,}000$; inconclusive). This reverse channel grows with $m$ on pgvector (size 0.21 to 0.48), which suggests a dependence
on the secrets, but pgvector's exact search shows an offset of the same sign (0.12 to 0.17, flat), no same-content comparison exists for constant work and the mechanism is not established, so the 98\% of pgvector's probes at
$m = 1{,}000$ that an attacker who knows the sign would identify is an upper bound on this harness. The script reads the rule's ``upper end at most 0.65'' at $m = 1{,}000$ as the window
$[0.40, 0.65]$ and prints these parts inconclusive; read one-sidedly, as worded, they would pass,
which would hide the channel.

\paragraph{Approximate pre-filtering (H3).}
pgvector's iterative scan separated the worlds as predicted (H3a confirmed): at $m = 1{,}000$ every probe was slower
on $W_1$ at $p = 10\%$ and 50\% (paired AUC 1.000; 13.5 against 4.7\,ms and 13.7 against 3.3\,ms), and so
was every probe at $m = 100$ and $p = 50\%$. The same scan returned \emph{better} answers on $W_1$:
recall@10 was 0.968 and 0.978 against 0.707 and 0.665 on $W_0^g$. This is consistent with how an iterative scan
works: made to pass the unreadable secrets, it ranks many more readable rows before it stops and so
comes close to exact search. An attacker thus learns of nearby secrets from the quality of its answer
as well as from its timing.
Elasticsearch (H3b, inconclusive) was slower on $W_1$ for most probes at $p = 10\%$ (0.800, 0.675 to
0.925), yet the mean latencies were 2.426 and 2.430\,ms, and the interval at $p = 50\%$ (0.463 to
0.750) straddles the threshold. Qdrant showed no channel in any observable at $p = 1\%$ and 10\% (H3c
confirmed; latency 0.460, 0.424 to 0.496), where its filter is small enough for an exact scan of the
payload index. At $p = 50\%$, outside the hypothesis, its filtered traversal had to cross the secrets:
recall@10 on $W_1$ fell to 0.530 against 0.708 (paired AUC 0.787, 0.688 to 0.875), and some answers
held fewer than ten hits, consistent with the island of Section~\ref{sec:qual-qdrant}, which we did not measure directly.

\paragraph{Exact pre-filtering (H4, inconclusive).}
Exact search returned the same answer in both worlds: count, completeness, recall and score gave a
paired AUC of 0.500 (0.500 to 0.500) on all four stores, pooled and at $m = 1{,}000$. Every H4 verdict
is nevertheless inconclusive, through latency alone. On Qdrant and Elasticsearch latency showed no
channel pooled (0.461 and 0.486), and their intervals at $m = 1{,}000$ (0.325 to 0.508 and 0.367 to
0.575) were too wide at 40 clusters for the script's window; under the rule's one-sided wording both
stores would be confirmed. On pgvector (0.147, 0.123 to 0.172) and FAISS (0.276, 0.246 to 0.306)
latency carried the build offsets described under H1, the same at every $m$ and so not about the
secrets. The follow-up traced pgvector's to the tables' state. At the start of every group the table
built last had no page marked all-visible, because autovacuum had not yet reached it (40 of 40
groups), while the table built at the start of the seed always had; in the main run the table built
last was $W_0^g$. Exact search slowed with the share of groups a table started unvacuumed: a mean of
59.9\,ms on $W_1$, never unvacuumed at a group's start, 61.4 and 74.3\,ms on later builds unvacuumed in
4 and 15 of the 40 groups, and 107\,ms on the build unvacuumed in 34. Once every table had been vacuumed, the four
ran the same plan with the same buffer counts in the same time. By the reading fixed in advance this
too is inconclusive (matched content 0.567, 0.530 to 0.602), and a small difference remains after
vacuuming that we cannot explain. Qdrant's exact search showed no difference between builds in any
comparison (0.493 to 0.505); the Elasticsearch part of the follow-up was lost when its runner stopped
responding, so Elasticsearch has no such comparison. The in-memory control failed G1.

\paragraph{Collection statistics (H5).}
Beside the results and prior work of Section~\ref{sec:side-results}: at $m = 1$ and 10 the AUCs were above the
falsification bound, so H5a is inconclusive for both stores. No cell isolates $m$
documents, at any $m$: $W_1$ and $W_0^g$ differ in all 1,233 secrets of the probe's group, which the zone rule keeps
beyond the reach of dense search but which collection statistics count. The shifts are small (a mean
score of 13.634 against 13.676 at $m = 1$ on Elasticsearch) and invisible to the one-shot attacker
(pooled AUC 0.503 to 0.531). \texttt{ts\_rank\_cd} showed no score channel and changed 0.2\% of hit
lists (H5b confirmed), although its latency, on which we made no prediction, separated the worlds at $m = 1{,}000$ (0.70, 0.58 to 0.81). Fusion changed the hit list for 56.7\% of probes at $m = 1{,}000$ (48.3\%
to 65.0\%; H5c confirmed), and for 26.9\% at $m = 1$.

\paragraph{Against prior work.}
Wang et al.'s STRESS attack writes documents to a shared index, maps its shards and uses score differences, or ranks alone, to discover terms in other tenants' documents and count the documents that hold a term; they propose document frequencies from public data as a countermeasure~\cite{wang2017sidechannel}. Espiritu and Cash extend this to indexed terms and, with search-as-you-type indexes, approximate text, through document-level security in Elasticsearch and OpenSearch, and recover records from PostgreSQL row-level security's timing with range and conjunctive predicates; in their tests cross-region placement and concurrent load lowered accuracy until more queries restored it, and a composite index with a subquery policy removed the timing channel for their scenario~\cite{espiritu2026plaintext}. Neither paper tests vector search. Our channels, measured in a forced-choice game with no inserted documents, are weaker claims: the score result re-measures their channel through the layer's pre-filter, and our latency channels come from the widening loop and the iterative scan. We did not measure row-level security's channels, although Appendix~\ref{sec:limits-native} reports its recall. The README of an open-source project reports count, relevance and timing channels of filter-after-retrieval RAG on a 35-document synthetic organization, in process, which we did not reproduce~\cite{raghuram2026onewayglass}.

\paragraph{Padding (H6).}
Seven cells had a latency channel before padding on seeds 1 to 4 (paired AUC 0.755 to 0.927): the
post-filter on FAISS and Qdrant at $m = 100$ and 1,000 and on Elasticsearch at 1,000, and pgvector's
pre-filter at 100 and 1,000. The analysis script computes each padded latency as $T \lceil x / T \rceil$, with $T$ the slowest latency of the first seed rounded up, so every recorded latency up to $T$ becomes $T$, and the padded AUC was
0.500 (0.500 to 0.500) in all seven by construction: both worlds get one value, and ties count one half. The script printed H6 inconclusive because its test required
five seeds of cells that by design have four; applied as written, the rule confirms H6, which adds nothing beyond that identity. The rule
covers only cells whose forward AUC reached 0.75. Reverse channels, such as constant work on pgvector
(paired AUC 0.017 at $m = 1{,}000$, $|\mathrm{AUC} - 0.5| = 0.48$), fell outside it and were not
tested, although a reverse channel serves an attacker equally.

\paragraph{Costs.}
Padding raised mean latency 5.2 to 6.7 times in those four configurations, from 1.33 to 27.1 times
over all configurations, and by medians of per-probe ratios up to 57.6, since every request waits at
least as long as the slowest answer of the first seed. Exact search cost 1.53, 2.01 and 8.10 times the mean latency of
the approximate pre-filter on Qdrant, Elasticsearch and pgvector (intervals 1.52 to 1.53, 1.99 to 2.03
and 7.73 to 8.48). Constant work cost 1.05 to 1.23 times the widening loop's mean on the three servers
and 0.74 on FAISS, where one round of 1,000 replaces up to four, but a median of 1.26 to 3.43 times
per probe, because it charges every cheap request the price of the expensive ones. A separated index
answered in 0.45 to 0.99 of the shared index's mean time in every configuration except pgvector's
keyword search (1.07, 0.98 to 1.18), as it holds fewer chunks and no secrets.

\paragraph{Summary.}
In this forced-choice game each enforcement path separated the worlds through the observable its
mechanism predicts: the widening loop through
latency once about 100 unreadable chunks lie near a query and through the hit count when the attacker
reads little, pgvector's approximate pre-filter through latency and recall, and keyword scores through
collection statistics. Latency channels smaller than the 0.08 to 0.38 that separated builds of
identical content, where we measured that, cannot be told from build noise unless they grow with $m$.
Constant work traded the widening channel for a reverse one on pgvector. On the shared index only
searches that ran exactly, whether requested or chosen by the planner (Qdrant's small filters), returned the same answer in both worlds, as exact search should, and a separated index does so by construction. Ideal padding costs five to seven times the mean latency of
the configurations that needed it; it was not tested against reverse channels, and the detection
fields of Appendix~\ref{sec:rep-design} remain an untested observable.
